\subsubsection{\protect\tr{Kernel} hữu tỉ}

Định nghĩa sau đây thiết lập một khuôn khổ tổng quát cho việc định nghĩa các
\tr{kernel} chuỗi.

\begin{defn}[\tr{Kernel} hữu tỉ]
  Một \tr{kernel} $K \colon \Sigma^\ast \times \Sigma^\ast \to \mathbb{R}$ được
  gọi là \tr{kernel} hữu tỉ (rational kernel) nếu nó trùng với một ánh xạ được
  định nghĩa bởi một \tr{transducer} có trọng số $U$: $\forall x, y \in
  \Sigma^\ast, K(x, y) = U(x, y)$.
\end{defn}

Lưu ý rằng ta hoàn toàn có thể chọn một định nghĩa tổng quát hơn: thay vì sử
dụng các \tr{transducer} có trọng số, ta có thể dùng những ánh xạ chuỗi mạnh
hơn, chẳng hạn như các phép chuyển đổi đại số (algebraic transductions)---là
đối tượng hàm tương ứng với các ngôn ngữ phi ngữ cảnh (context-free
languages)---hoặc thậm chí những mô hình còn mạnh hơn nữa.

Tuy nhiên, một yêu cầu thiết yếu đối với \tr{kernel} là khả năng tính toán hiệu
quả, và các định nghĩa phức tạp hơn sẽ dẫn đến độ phức tạp tính toán của việc
tính \tr{kernel} tăng lên đáng kể.

Đối với \tr{kernel} hữu tỉ, tồn tại một thuật toán tính toán tổng quát và hiệu
quả.

\paragraph{Tính toán}

Chúng ta sẽ giả định rằng \tr{transducer} $U$ dùng để định nghĩa một
\tr{kernel} hữu tỉ $K$ không chứa chu trình $\epsilon$ ($\epsilon$-cycle) nào
có trọng số khác không; nếu không, giá trị của hàm \tr{kernel} sẽ là vô hạn đối
với mọi cặp chuỗi. Với chuỗi $x$ bất kì, gọi $T_x$ là một \tr{transducer} có
trọng số chỉ có duy nhất một đường đi chấp nhận, trong đó cả nhãn đầu vào và
nhãn đầu ra đều là $x$, và trọng số của nó bằng 1. $T_x$ có thể được xây dựng
một cách dễ dàng từ $x$ trong thời gian tuyến tính $O\bigl( |x| \bigr)$. Khi
đó, với mọi $x, y \in \Sigma^*$, giá trị $U(x, y)$ có thể được tính toán thông
qua hai bước sau: \begin{enumerate}
  \item Tính $V = T_x \circ U \circ T_y$ bằng cách sử dụng thuật toán hợp với
    độ phức tạp thời gian là $O\bigl(|U||T_x||T_y|\bigr)$.
  \item Tính tổng trọng số của tất cả các đường đi chấp nhận của $V$ bằng cách
    sử dụng thuật toán tìm khoảng cách ngắn nhất tổng quát (general
    shortest-distance algorithm) với độ phức tạp thời gian là $O\bigl( |V|
    \bigr)$.
\end{enumerate}
Theo định nghĩa của phép hợp, $V$ là một \tr{transducer} có trọng số mà các
đường đi chấp nhận của nó chính là những đường đi chấp nhận của $U$ có nhãn đầu
vào là $x$ và nhãn đầu ra là $y$. Bước thứ hai sẽ tính tổng trọng số của các
đường đi này, tức là bằng chính xác $U(x, y)$. Vì $U$ không chứa chu trình
$\epsilon$ nên $V$ không có chu trình, và bước này có thể được thực hiện trong
thời gian tuyến tính. Do đó, độ phức tạp tổng thể của thuật toán để tính $U(x,
y)$ là $O\bigl(|U||T_x||T_y|\bigr)$. Vì $U$ là cố định đối với \tr{kernel} hữu
tỉ $K$, và $|T_x| = O\bigl(|x|\bigr)$ với mọi $x$, điều này cho thấy các giá
trị của hàm \tr{kernel} có thể tính được trong thời gian bậc hai $O\bigl(|x|
|y| \bigr)$. Đối với một số \tr{transducer} có trọng số $U$ cụ thể, việc tính
toán có thể hiệu quả hơn, ví dụ như trong $O(|x| + |y|)$.

\paragraph{\protect\tr{Kernel} hữu tỉ PDS}

Với bất kỳ \tr{transducer} $T$ nào, gọi $T^{-1}$ là bộ chuyển đổi nghịch đảo
của $T$, tức là \tr{transducer} thu được từ $T$ bằng cách hoán đổi nhãn đầu vào
và nhãn đầu ra của mọi \tr{transition}. Với mọi $x, y$, ta luôn có $T^{-1}(x,
y) = T(y, x)$. Định lý sau đây đưa ra một phương pháp tổng quát để xây dựng một
\tr{kernel} hữu tỉ PDS từ một \tr{transducer} có trọng số bất kì.

\begin{thm}
  Với mọi \tr{transducer} có trọng số $T = (\Sigma, \Delta, I, F, E, \rho)$,
  hàm $K = T \circ T^{-1}$ là một \tr{kernel} hữu tỉ PDS.
\end{thm}

\begin{proof}
  Theo định nghĩa của phép hợp và nghịch đảo, với mọi $x, y \in \Sigma^\ast$,
  \begin{equation}
    K(x, y) = \sum_{z \in \Delta^\ast} T(x, z) T(y, z).
  \end{equation}$K$ là giới hạn điểm của dãy \tr{kernel} $(K_n)_{n\ge 0}$ định
  nghĩa bởi: \begin{equation}
    \forall n \in \mathbb{N}, \forall x, y \in \Sigma^\ast, \quad K_n(x, y) =
    \sum_{|z| \le n} T(x, z) T(y, z),
  \end{equation} với tổng chạy trên tất cả chuỗi trong $\Delta^\ast$ có độ dài
  không quá $n$. $K_n$ là PDS vì ma trận \tr{kernel} nhân tương ứng
  $\mathbf{K}_n$ với mọi mẫu $\{x_1, \dots, x_m\}$ là SPSD. Điều này dễ thấy vì
  $\mathbf{K}_n$ có thể viết thành $\mathbf{K}_n = \mathbf{A}
  \mathbf{A}^\intercal$ với $\mathbf{A} = \bigl( K_n(x_i, z_j) \bigr )_{i \in
  [m], j \in [N]}$, trong đó $z_1,\dots, z_N$ là phép liệt kê tùy ý trong tập
  các chuỗi kí tự trong $\Sigma^\ast$ với độ dài không quá $n$. Do đó, $K$ là
  PDS vì nó là giới hạn điểm của dãy các \tr{kernel} PDS $(K_n)_{n \in
  \mathbb{N}}$.
\end{proof}

Các \tr{kernel} chuỗi thường được dùng trong sinh học tính toán (computational
biology), xử lí ngôn ngữ tự nhiên, thị giác máy tính và các ứng dụng khác là
tất cả trường hợp đặc biệt của \tr{kernel} hữu tỉ có dạng $T \circ T^{-1}$. Các
\tr{kernel} có thể được tính hiệu quả thông qua cùng thuật toán tổng quát để
tính \tr{kernel} hữu tỉ được trình bày ở đoạn trước. Bởi vì \tr{transducer} $U
= T \circ T^{-1}$ định nghĩa \tr{kernel} hữu tỉ PDS có dạng cụ thể, có nhiều
lựa chọn khác nhau để tính $T_x \circ U \circ T_y$: \begin{itemize}
  \item tính $U = T \circ T^{-1}$ trước sau đó đến $V = T_x \circ U \circ T_y$;
  \item tính $V_1 = T_x \circ T$ và $V_2 = T_y \circ T$ trước, sau đó tính $V =
    V_1 \circ V_2^{-1}$;
  \item tính trước $V_1 = T_x \circ T$1, sau đó $V_2 = V_1 \circ T^{-1}$, sau
    đó $V = V_2 \circ T_y$, hoặc chuỗi các phép toán tương tự với $x$ và $y$
    được hoán vị.
\end{itemize} Các phương pháp trên đều có cùng kết quả khi tính tổng các
trọng số của đường đi chấp nhận và có cùng độ phức tạp trong trường hợp xáu
nhất. Tuy nhiên, trong thực tế, do độ thưa các kết quả hợp trung gian, có thể
có sự khác biệt lớn về chi phí tính toán về mặt thời gian và không gian giữa
chúng. Một phương pháp thay thế dựa trên hợp $n$ thành phần ($n$-way
composition) có thể giúp tính toán hiệu quả hơn.

\begin{exmp}[\tr{Kernel} chuỗi \tr{bigram} và \tr{gappy bigram}]
  \figurename~\ref{fig:bigram-transducer} cho thấy một \tr{transducer} có trọng
  số $T_\text{bigram}$ định nghĩa một \tr{kernel} chuỗi phổ phổ biến,
  \emph{\tr{kernel} chuỗi bigram} (bigram sequence kernel), cho trường hợp cụ
  thể của bảng chữ cái được thu hẹp còn $\Sigma = \{a, b\}$. \tr{Kernel}
  \tr{bigram} gắn với với chuỗi $x$ và $y$ bất kì tổng các tích của đếm tất cả
  \tr{bigram} trong $x$ và $y$. Với mọi chuỗi $x \in \Sigma^\ast$ và
  \tr{bigram} $z \in \{aa, ab, ba, bb\}$, $T_\text{bigram}(x, z)$ chính xác là
  số lần xuất của \tr{bigram} $z$ trong $x$. Do đó, theo định nghĩa của phép
  hợp, và nghịch đảo, $T_\text{bigram} \circ T_\text{bigram}^{-1}$ tính chính
  xác \tr{kernel} \tr{bigram}.

  \begin{figure}[htbp]
    \centering

    \begin{subfigure}{\textwidth}
      \centering
      \import{\figurespath}{bigram-transducer}
      \caption{}
      \label{fig:bigram-transducer}
    \end{subfigure}

    \begin{subfigure}{\textwidth}
      \centering
      \import{\figurespath}{gappy-bigram-transducer}
      \caption{}
      \label{fig:gappy-bigram-transducer}
    \end{subfigure}
    \caption[\protect\tr{Transducer} \protect\tr{bigram} và \protect\tr{gappy
    bigram}]{\subref*{fig:bigram-transducer}~\tr{Transducer} $T_\text{bigram}$
    định nghĩa \tr{kernel} \tr{bigram} $T_\text{bigram} \circ
    T_\text{bigram}^{-1}$ với $\Sigma = \{a, b\}$.
    \subref*{fig:gappy-bigram-transducer}~\tr{Transducer}
    $T_\text{gappy\_bigram}$ định nghĩa \tr{kernel} \tr{gappy bigram}
    $T_\text{gappy\_bigram} \circ T_\text{gappy\_bigram}^{-1}$ với phạt khoảng
    trổng $\lambda \in (0, 1)$.}
  \end{figure}

  \figurename~\ref{fig:gappy-bigram-transducer} cho thấy một \tr{transducer} có
  trọng số $T_\text{gappy\_bigram}$ định nghĩa \emph{\tr{kernel} \tr{gappy
  bigram}} (gappy bigram kernel). \tr{Kernel} \tr{gappy bigram} gắn với hai
  chuỗi $x$ và $y$ tổng các tích của đếm tất cả \tr{gappy bigram}
  trong $x$ và $y$ và phạt bởi độ dài khoảng trống (gap) của chúng. \tr{Gappy
  bigram} là những chuỗi có dạng $aua$, $aub$, $bua$, $bub$, với $u \in
  \Sigma^\ast$ được gọi là khoảng trống. Kết quả đếm của một \tr{gappy bigram}
  được nhân với $\lambda^{|u|}$ với $\lambda \in (0, 1)$ cố định, do đó
  \tr{gappy bigram} với khoảng trống càng dài thì đóng góp càng ít vào định
  nghĩa của sự tương đồng. Trong khi định nghĩa này có vẻ phức tạp,
  \figurename~\ref{fig:gappy-bigram-transducer} cho thấy
  $T_\text{gappy\_bigram}$ có thể suy ra trực tiếp từ $T_\text{bigram}$. Việc
  biểu diễn \tr{kernel} hữu tỉ bằng đồ thị giúp ta hiểu hoặc tùy biến định
  nghĩa của chúng trực quan hơn.
\end{exmp}

\paragraph{\protect\tr{Transducer} đếm}

Định nghĩa của hầu hết các \tr{kernel} chuỗi đều dựa trên số lần xuất hiện (số
đếm) của một số mẫu (patterns) phổ biến trong các chuỗi dữ liệu. Trong các ví
dụ vừa được khảo sát ở trên, các mẫu này chính là các bigram hoặc các \tr{gappy
bigram}. Hiện tồn tại một phương pháp đơn giản và tổng quát để xây dựng một
\tr{transducer} có trọng số nhằm đếm số lần xuất hiện của các mẫu, từ đó sử
dụng chúng để định nghĩa các \tr{kernel} hữu tỉ PDS. Gọi $\mathcal{X}$ là một
\tr{automaton} hữu hạn biểu diễn tập hợp các mẫu cần đếm. Trong trường
hợp của các hạt nhân bigram với bảng chữ cái $\Sigma = \{a, b\}$, $\mathcal X$
sẽ là một \tr{automaton} chấp nhận chính xác tập hợp các chuỗi $\{aa, ab, ba,
bb\}$.  Khi đó, \tr{transducer} có trọng số trong
\figurename~\ref{fig:counting-transducer} có thể được sử dụng để tính toán
chính xác số lần xuất hiện của từng mẫu được chấp nhận bởi $\mathcal X$.

\begin{figure}[htbp]
  \centering
  \import{\figurespath}{counting-transducer}
  \caption[\protect\tr{Transducer} đếm]{\tr{Transducer} đếm $T_\text{count}$
  với $\Sigma = \{a, b\}$. ``\tr{transition}'' X:X/1 đại diện cho một
  \tr{transducer} có trọng số được tạo từ một \tr{automaton} $\mathcal{X}$ bằng
  cách thêm vào mỗi \tr{transition} một nhãn đầu ra đồng nhất với nhãn hiện có,
  và gán tất cả trọng số ở mỗi \tr{transition} và trọng số kết thúc bằng 1.}
  \label{fig:counting-transducer}
\end{figure}

\begin{thm}
  Với mọi $x \in \Sigma^\ast$ và chuỗi $z$ được chấp nhận bởi $\mathcal{X}$,
  $T_\text{count}(x, z)$ là số lần xuất hiện của $z$ trong $x$.
\end{thm}

\begin{proof}
  Gọi $x \in \Sigma^\ast$ là một chuỗi bất bất kì và $z$ là một chuỗi được chấp
  nhận bởi $\mathcal{X}$. Bởi vì mọi đường đi chấp nhận của $T_\text{count}$
  đều có trọng số là 1 nên $T_\text{count}(x, z)$ bằng với số đường đi chấp
  nhận trong $T_\text{count}$ nhãn đầu vào là $x$ và nhãn đầu ra $z$.

  Để ý rằng, mỗi đường đi chấp nhận $\pi$ trong $T_\text{count}$ có thể phân
  tích thành $\pi = \pi_0\pi_{01}\pi_1$, trong đó $\pi_0$ là đường đi chỉ qua
  vòng của trạng thái 0 với nhãn đầu vào là tiền tố $x_0$ và nhãn đầu ra
  $\epsilon$, $\pi_{01}$ là đường đi chấp nhận từ 0 đến 1 với nhãn đầu vào và
  đầu ra đều là $z$, và $\pi_1$ là đường đi chỉ qua vòng tại trạng thái 1 với
  nhãn đầu vào là hậu tố $x_1$ và nhãn đầu ra là $\epsilon$. Do đó, số đường đi
  chính là số cách khác nhau biểu diễn $x$ thành $x=x_0zx_1$, tức là số lần
  xuất hiện của $z$ trong $x$.
\end{proof}

Định lý trên cung cấp một phương pháp rất tổng quát để xây dựng các \tr{kernel}
hữu tỉ PDS dạng $T_{\text{count}} \circ T_{\text{count}}^{-1}$ dựa trên số lần
xuất hiện (số đếm) của các mẫu có thể được định nghĩa thông qua một
\tr{automaton} hữu hạn, hoặc tương đương là một biểu thức chính quy (regular
expression). \figurename~\ref{fig:counting-transducer} hiển thị bộ chuyển đổi
cho trường hợp bảng chữ cái đầu vào được thu hẹp còn $\Sigma = \{a, b\}$.
Trường hợp tổng quát có thể được suy ra một cách trực tiếp bằng cách bổ sung
thêm các vòng vào các trạng thái 0 và 1 bằng các ký tự khác ngoài $a$ và
$b$. Trong thực tế, cơ chế tính toán lười (lazy evaluation) có thể được áp dụng
để tránh việc phải khởi tạo tường minh các \tr{transition} này cho toàn bộ các
ký tự trong bảng chữ cái; thay vào đó, chúng chỉ được tạo ra theo nhu cầu  dựa
trên các ký tự thực tế xuất hiện trong chuỗi đầu vào $x$. Cuối cùng, ta
có thể gán các trọng số khác nhau cho các mẫu được đếm để nhấn mạnh hoặc giảm
bớt tầm quan trọng của một số mẫu nhất định, giống như trong trường hợp của các
\tr{gappy bigram}. Điều này có thể được thực hiện một cách đơn giản bằng cách
thay đổi trọng số của các bước chuyển hoặc trọng số kết thúc của \tr{automaton}
$\mathcal X$ được dùng để định nghĩa $T_{\text{count}}$.
